\chapter{Введение}
\label{ch:intro}

    В ряде задач для производства алюминиевого проката необходимо определить оптимальную мощность печи.  
    
    Можно опытным путём попытаться найти наиболее оптимальную мощность печи, однако большое количество материалов будет израсходовано на данные опыты, а также найденная мощность будет привязана к одному размеру порции алюминия, а для её корректировки вследствие изменения рынка алюминиевых деталей потребуется заново проводить данные опыты, что повлечёт за собой большие расходы. Следовательно, этот способ не является эффективным.
    
    В качестве второго варианта выступает расчёт энергопотребления. Рассчитать необходимую мощность можно заранее, используя теоретические сведения об алюминии. Этот вариант подразумевает меньшие расходы, поэтому является наиболее подходящим. Перед определением требуемой мощности необходимо рассчитать, какое количество теплоты поглощается алюминием при нагревании на 1 Кельвин. Данная величина называется теплоёмкостью. 

    Конечной целью работы является определение удельной теплоёмкости твёрдого тела.

    % + для различных сплавов теплоёмкость может отличаться от справочных значений, поэтому следует взять небольшую порцию алюминия и измерить 